\ifx\allfiles\undefined
\documentclass[12pt, a4paper, oneside, UTF8]{ctexbook}
\def\path{../config}
\input{\path/package.tex}

% 定理环境设置
\input{\path/theorem1_zh.tex}

\input{\path/custom.tex}

% 修改参数改变封面样式，0 默认原始封面、内置其他1、2、3种封面样式
\def\myIndex{3}


\ifnum\myIndex>0
    \input{\path/cover_package_\myIndex} 
\fi

\def\myTitle{高等数学笔记}
\def\myAuthor{M.K.}
\def\myDateCover{\today}
\def\myDateForeword{\today}
\def\myForeword{前言}
\def\myForewordText{
    
    这是一个基于\LaTeX{}模板的高数笔记．
    内容参考了包括但不限于以下资料：
    \begin{itemize}
        \item 《高等数学》（同济大学出版社）第七版
        \item 《吉米多维奇高等数学习题精选精讲》（山东科学技术出版社）第二版
        \item 《高等数学作业集》（哈尔滨工业大学出版社）第一版
        \item \href{https://space.bilibili.com/1035929235}{B站UP主：一高数}的高数课程
        \item \href{https://space.bilibili.com/42428180}{B站UP主：考研竞赛数学}的试题讲解与高数课程
        \item \href{https://space.bilibili.com/391930545}{B站UP主：Maki的完美算数教室}的高数内容

    \end{itemize}

    封面来源于\href{https://www.pixiv.net/users/3952}{おにねこ}的作品\href{https://www.pixiv.net/artworks/136551599}{雨雫とプレアデス}
    如有侵权请联系我删除．
}
\def\mySubheading{副标题}


\begin{document}
% \input{../config/cover}
\else
\fi
\chapter{微分中值定理的证明题}
\section{微分中值定理的证明题练习}
\begin{exercise}
    设函数$f(x)$在区间$[0,1]$上可导，且$f(0)>0,f(1)>0,\int_{0}^{1}f(x) \d x = 0$。证明：
    \begin{enumerate}[label=(\arabic*)]
        \item $f(x)$在区间$[0,1]$内至少有两个不同的零点；
        \item $\exists \xi \in (0,1)$使得$f'(\xi)+3f^{3}(\xi)=0$
    \end{enumerate}
\end{exercise}
\begin{proof}
    \begin{enumerate}[label=(\arabic*)]
        \item 若$f(x)$在区间$[0,1]$内恒大于零或恒小于零，则$\int_{0}^{1}f(x) \d x$应大于零或小于零，与题意矛盾。故至少$\exists x_0,$使得$f(x_0) <0$.
        又因为$f(0)>0,f(1)>0$，所以$f(x)$在区间$(0,x_0)$与$(x_0,1)$内至少各有一个零点，故$f(x)$在区间$[0,1]$内至少有两个不同的零点。
        \item 设$f(x)$在区间$[0,1]$内的两个不同零点为$x_1,x_2$，令$F(x)=f(x)e^{\int_{0}^{x}3f^{2}(t) \d t}$，则
        \[
        F'(x)=(f'(x)+3f^{3}(x))e^{\int_{0}^{x}3f^{2}(t) \d t}, \quad F(x_1)=F(x_2)=0
        \]
        由罗尔定理知，$\exists \xi \in (x_1,x_2) \subset (0,1)$使得$F'(\xi)=0$，即$f'(\xi)+3f^{3}(\xi)=0$。
    \end{enumerate}
\end{proof}
\begin{exercise}
    已知$f(x)=0,f''(x)$存在，证明：\\$\text{\qquad \qquad \qquad \qquad}$
    $\exists \xi \in (-\dfrac{\pi}{2},\dfrac{\pi}{2})$，使得$f''(\xi)-3\tan \xi \cdot f'(\xi)-2f(\xi)=0$
\end{exercise}
\begin{proof}
    令$F(x)=f(x)\cos^2 x$，则
    \[
        F'(x)=f'(x)\cos^2 x - 2f(x)\cos x \sin x
    \]
    \[
        F''(x)=f''(x)\cos^2 x - 4f'(x)\cos x \sin x - 2f(x)(\cos^2 x - \sin^2 x)
    \]
    因为$F(0)=F(-\dfrac{\pi}{2})=F(\dfrac{\pi}{2})=0$，由罗尔定理知，$\exists \xi_1 \in (-\dfrac{\pi}{2},0), \xi_2 \in (0, \dfrac{\pi}{2})$使得$F'(\xi_1)=F'(\xi_2)=0$\\
    再次使用罗尔定理，$\exists \xi \in (\xi_1, \xi_2) \subset (-\dfrac{\pi}{2},\dfrac{\pi}{2})$使得$F''(\xi)=0$，即证。
\end{proof}
\begin{myRemark}
    万能构造：$f'(\xi)+g(\xi)f(\xi)=0 \Rightarrow F(x)=f(x)e^{\int g(t) \d t}$
\end{myRemark}
\begin{exercise}
    设函数$f(x)$在区间$[a,b]$上连续,在区间$(a,b)$内可导，且$\exists c \in (a,b)$使得$f'(c)=0$。证明：
    $\exists \xi \in (a,b)$使得$f'(\xi)=\dfrac{f(\xi)-f(a)}{b-a}$
\end{exercise}
\begin{proof}
    令$F(x)=[f(x)-f(a)]e^{-\dfrac{x}{b-a}}$，则
    \[
        F'(x)=\left[f'(x)-\dfrac{f(x)-f(a)}{b-a}\right]e^{-\dfrac{x}{b-a}}, \quad F(c)= \left[f(c)-f(a)\right]e^{-\dfrac{c}{b-a}}, \quad F(a)=0
    \]
    \begin{enumerate}[label=\roman*).]
        \item 若$f(c)=f(a)$，则$F(c)=0$，故由罗尔定理知，$\exists \xi \in (a,c) \subset (a,b)$使得$F'(\xi)=0$，即$f'(\xi)=\dfrac{f(\xi)-f(a)}{b-a}$；
        \item 若$f(c)\neq f(a)$，不妨设$f(c) > f(a)$，则
        \[
            F'(c) = -\dfrac{f(c)-f(a)}{b-a}e^{-\dfrac{c}{b-a}} < 0
        \]
        由拉格朗日中值定理知，$\exists \eta \in (a,c)$使得
        \[
            F'(\eta) = \dfrac{F(c)-F(a)}{c-a} > 0
        \]
        故由介值定理可知，$\exists \xi \in (\eta,c) \subset (a,b)$使得$F'(\xi)=0$，即$f'(\xi)=\dfrac{f(\xi)-f(a)}{b-a}$
    \end{enumerate}
\end{proof}
\begin{exercise}
    设函数$f(x)$在区间$[0,1]$内连续，在区间$(0,1)$内可导，且$f(0)=1,f(1)=\dfrac{1}{2}$．证明：$\exists \xi \in (0,1)$使得$f'(\xi) + f^2 (\xi) =0$
\end{exercise}
\begin{proof}
    \begin{enumerate}[label=\roman*).]
        \item 若$f(x)$无零点，则令$F(x)=-\dfrac{1}{f(x)}+x$，则
        \[
           F(0)=F(1)=-1 , \quad F'(x)=\dfrac{f'(x)}{f^2 (x)} + 1
        \]
        由罗尔定理知，$\exists \xi \in (0,1)$使得$F'(\xi)=0$，即$f'(\xi) + f^2 (\xi) =0$；
        \item 若$f(x)$有不止一个零点，不妨设$f(x_1)=f(x_2)=\cdots = f(x_n)=0$，则令$F(x)=f(x)e^{\int_{0}^{x} f(t) \d t}$，则
        \[
            F'(x)=(f'(x)+f^2 (x))e^{\int_{0}^{x} f(t) \d t}, \quad F(x_1)= F(x_2)=\cdots = F(x_n)=0
        \]
        由罗尔定理知，$\exists \xi \in (x_1,x_n) \subset (0,1)$使得$F'(\xi)=0$，即$f'(\xi) + f^2 (\xi) =0$；
        \item 若$f(x)$有唯一零点，不妨设$f(x_0)=0$\\
        若$x_0$并非$f(x)$在区间$(0,1)$上的最小值点，则$\exists r \in (0,1)$使得$f(r) < 0$，又因为$f(1) >0,f(0) >0$，故$f(x)$至少在区间$(0,r)$与$(r,1)$内各有一个零点，矛盾；\\
        因此,$x_0$为$f(x)$在区间$(0,1)$上的最小值点，故$f'(x_0)=0$，即$f'(x_0) + f^2 (x_0) =0$。
    \end{enumerate}
\end{proof}
\begin{myRemark}
    \[
        f'(\xi) + g(\xi)f^n (\xi) =0 \Rightarrow F(x)=\int \dfrac{f'(x)}{f^n (x)} \d x + \int g(x) \d x
    \]
\end{myRemark}
\begin{exercise}
    设函数$f(x)$在区间$[0,3]$上二阶可导，且满足$2f(0)=\int_{0}^{2}f(x)\d x=f(2)+f(3)$。证明：$\exists \xi \in (0,3)$使得$f''(\xi)-2f'(\xi)=0$
\end{exercise}
\begin{proof}
    令$F(x)=f'(x)e^{-2x}$，则
    \[
        F'(x)=e^{-2x}\left[f''(x)-2f'(x)\right]
    \]
    由$f(0)=\dfrac{1}{2}\displaystyle\int_{0}^{2}f(x)\d x=\dfrac{f(2)+f(3)}{2}$，
    可得$\exists x_1 \in (0,2),x_2 \in (2,3)$使得$f(0)=f(x_1)=f(x_2)$，
    由罗尔定理知，$\exists \xi_1 \in (0,x_1),\xi_2 \in (x_1,x_2)$使得$f'(\xi_1)=f'(\xi_2)=0$，即$F(\xi_1)=F(\xi_2)=0$\\
    再次使用罗尔定理，$\exists \xi \in (\xi_1,\xi_2) \subset (0,3)$使得$F'(\xi)=0$，即$f''(\xi)-2f'(\xi)=0$。
\end{proof}
\begin{exercise}
    设函数$f(x)$在区间$[-2,2]$上二阶可导，且$|f(x)| \le 1,[f(0)]^2 +[f'(0)]^2=4$,证明：$\exists \xi \in (-2,2)$使得$f''(\xi)+f(\xi)=0$
\end{exercise}
\begin{proof}
    令$F(x)=[f(x)]^2 +[f'(x)]^2$，则
    \[
        F'(x)=2f(x)f'(x)+2f'(x)f''(x)=2f'(x)[f(x)+f''(x)]
    \]
    由拉格朗日中值定理可得
    \begin{gather*}
        f(-2)-f(0)=-2f'(x_1),x_1 \in (-2,0)\\
        f(2)-f(0)=2f'(x_2),x_2 \in (0,2)
    \end{gather*}
    于是有$|f'(x_1)|=\dfrac{|f(-2)-f(0)|}{2} \le \dfrac{|f(2)|+|f(0)|}{2} < 1$,\quad 同理有$|f'(x_2)|<1$\\
    故$F(x_1) < 1^2 + 1^2 = 2<f(0)$,\quad 同理有 $F(x_2) < 1^2 + 1^2 = 2<f(0)$\\
    因此$f(x)$在区间$[-2,2]$上的最大值点$x_0$一定有$x_0 \in (x_1,x_2) $，则$f'(x_0)=0\Rightarrow f'(x_0)[f(x_0)+f''(x_0)]=0$\\
    若$f(x_0)=0$，则$F(x_0)=[f(x_0)]^2 +[f'(x_0)]^2 = [f'(x_0)]^2 <1<4=F(0)$，矛盾，故$f(x_0) \neq 0$，则$f(x_0)+f''(x_0)=0$，即存在$\xi \in (-2,2)$使得$f''(\xi)+f(\xi)=0$。
\end{proof}
\begin{exercise}
    设函数$f(x)$在区间$[0,1]$上连续，在区间$(0,1)$内可导，且$f(0)=0, \forall x \in [0,1],f(x) >0.$证明：$\forall t >0,\exists \xi \in (0,1)$使得$\dfrac{tf'(\xi)}{f(\xi)}=\dfrac{f'(1-\xi)}{f(1-\xi)}$
\end{exercise}
\begin{proof}
    令$F(x)=f(x)e^{\displaystyle\int -\dfrac{f'(1-x)}{t f(1-x)} \d x}=f(x)\left[f(1-x)\right]^{\tfrac{1}{t}}$\\
    则$F(0)=F(1)=0$，由罗尔定理知，$\exists \xi \in (0,1)$使得$F'(\xi)=0$，即$\dfrac{tf'(\xi)}{f(\xi)}=\dfrac{f'(1-\xi)}{f(1-\xi)}$。
\end{proof}
\begin{exercise}
    设函数$f(x)$在区间$[0,1]$上二阶可导，且$f(0)=f'(0)=0.$证明：$\exists \xi \in (0,1)$使得$f''(\xi)=\dfrac{2f(\xi)}{(1-\xi)^2}$
\end{exercise}
\begin{proof}
    $f''(x)(1-x)^2 - 2f(x) = [f'(x)(1-x)^2 - 2f(x)(1-x)]'$\\
    令$F(x)=f'(x)(1-x)^2 - 2f(x)(1-x)$，则
    \[
        F(0)=f'(0)-2f(0)=0, \quad F(1)=0
    \]
    由罗尔定理知，$\exists \xi \in (0,1)$使得$F'(\xi)=0$，即$f''(\xi)=\dfrac{2f(\xi)}{(1-\xi)^2}$.
\end{proof}
\begin{myRemark}
    \[
        \dfrac{f(x)}{g(x)} = \dfrac{f''(x)}{g''(x) }\Rightarrow [f(x)g'(x)-f'(x)g(x)]'=0
    \]
\end{myRemark}
\ifx\allfiles\undefined
\end{document}
\fi